\section[第七章\q 资金管理]{第七章\\ 资金管理}

\subsection{系统与资金管理}

{\kaishu
资金管理只能将交易系统产生的劣势进行转换，而根本无法将其消除。所以说，任何资金管理策略都无法打败概率。

\hfill ——金融帝国
}

中国以前很流行赌骨牌。骨牌共有36个数字，从1到36。赌客可以任意对数字下注，庄家开牌时只开出一个数字。如果被你压中的话，1赔35。这时就出现了一条“生财之道”，赌客可以对其中34个数字下注。只要庄家没有开出赌客遗漏的那两个数字，那么赌客就能够获得下注额3％的收入。我们不难计算，这种下注方式的正确率为94.4％，而赔率为34赔1。这时赌客就可以总是赢钱，但千万不能碰上那少数失败的案例。

大家是否有一种似曾相识的感觉呢？事实上，大部分股民与上面的赌客并没有任何的不同。非常大的正确率优势与非常大的赔率劣势，组合成的一种具有负期望收益的交易策略。事实上，这并非生财之道，反而是一种自欺欺人的手段。靠必然出现的偶然事件走向破灭的交易者，只不过是在进行“零取整存”的无聊游戏。遗憾的是，有太多的人都对这种无聊游戏乐此不疲！

事实上，赌博是靠两部分来实现的。一部分是游戏规则，另一部分是下注策略。对于交易而言就是，一部分是交易系统，另一部分是资金管理。单纯的交易系统能够决定你交易时产生的期望收益，而资金管理则是你的下注策略。

按照上面的案例，如果说游戏规则决定了期望收益优势偏向于庄家的话，那么赌客无论采用什么样的下注策略都是无法打败庄家的。也就是说，如果交易者的交易系统并不存在期望收益优势的话，那么资金管理的优劣只能决定交易者破产速度的快慢。

资金管理只能将交易系统产生的劣势进行转换，而根本无法将其消除。这里我要说的就是加倍下注法。如果下注策略单次失败率为5.6％，那么连续失败两次的概率就为0.31％；而连续失败三次的概率就是0.0176％。这时我们就可以在少数出现失败的情况下，将下注额提高34倍。如果再次出现失败的话，那么就将下注额继续提高34倍。从概率上来讲，你赌一万次才可能碰上一次连续输三次的倒霉事！但反过来讲，你为了赢得一元钱的赌博收益，要准备3.93万元的后续资金。即使这样，你长期赌博下去，早晚也还是会遇上将3.93万元输光的结局。如果在市场上你准备用1万元来赚取1元利润的话，那么任何人都可以给你设计出非常安全的下注模式。问题是，没有人愿意动用1万元来赢得1元的投资收益，所以更多的人都是在没有足够后续资金的情况下使用加倍下注法。

那么是否说明资金管理策略是完全没有意义的呢？不，绝对不是这样！事实上，很多赌场都是在具有期望收益优势的前提下，仍旧采用严格的资金管理策略。通常来说，赌场都有最高下注额的限制。如果抛开这种限制的话，那么照样有打败赌场的方法。如果赌场的全部资金只有1个亿，而你拿出34个亿的资金在34个数字上下注的话，那么就有94.4％的概率能够将赌场彻底打败。事实上，赌场相对于赌客的优势，还包括拥有更多资金。

资金管理策略并不能创造出财富，但它又是交易者不可或缺的“减震器”。即使你的交易系统具有明显的期望收益，但只要你的下注足够重，那么市场照样可以将你扫地出门。虽然我们可以通过期望收益在宏观上计算出多次下注后获得收益的理论值，但在微观上的几次真实操作后，战绩仍旧可能出现相对于理论值不小的偏离。只要这种偏离理论值的幅度能够超出资金管理策略的承受能力，那么顷刻间你的全部交易行为就会出现崩溃。从某种意义上来说，资金管理策略的稳定性要比交易系统的期望收益更加重要。

找到仓位大小的最佳平衡点是一件非常困难的事情，这需要对交易系统进行严格的测试。我并不提倡找到理论上最佳的资金管理策略，相反我提倡用保守性原则来选择资金管理策略。很显然，如果你的资金管理策略需要使用复杂的计算，那么你所承受的风险毫无疑问是太大了。何况，我们脆弱的心灵是否能够承受过大的资金回撤幅度，也是一个显而易见的问题。要知道，如果轻仓下你都无法实现稳定获利，那么重仓之下就一定不可能实现稳定获利。反之，即使轻仓下你能够实现稳定获利，那么重仓下也不能保证仍旧能够实现稳定获利。

当我们在阅读交易书籍时会发现，几乎所有的作者都会强调资金管理的重要性，而又很少有人能够详细讨论资金管理的问题。我不得不承认，资金管理的章节应该是最难写的部分。出现这种情况的原因，并非是因为资金管理的问题有多么深奥，而只是因为资金管理是一个“二级问题”。回答这样的问题，就好像回答什么样的制动方式才是最好的。很显然，这个问题的本身并不复杂，但在不知道具体是什么样的交通工具之前，回答这个问题是没有意义的。汽车、飞机、轮船、自行车，不同交通工具使用的制动方式显然是不同的。既然脱离了交易系统来讨论资金管理是没有意义的，那么我们对于资金管理的讨论，或许只能够停留在一些原则性的问题上。

其实，即使是一些关于资金管理的原则性问题，对于不同的交易者而言，答案都有可能是不同的。比方说，“赢冲输缩”就是一种大部分交易者都认同的原则。但这种原则运用在不同的交易系统上，很可能出现不同的效果。甚至，实现“赢冲输缩”的具体方法，也是可以从系统开发而不是资金管理上来实现。

我始终觉得，既然交易系统是按照某种“一致性”的机械策略来应对市场，那么系统输出的盈亏分布就不太可能是杂乱无章的随机分布。对于具有正期望收益，并且结果是随机分布的赌博方式来说，“赢冲输缩”的策略无疑是正确的。但如果盈亏交易的分布具有某种明显的相关性，那么“赢冲输缩”的策略就要被我们重新审视一下了。这里我们可以举一个简单的例子：如果你在押大小的赌博中连续输了十次，那么第十一次赌博的结果仍旧是完全不具备任何概率上的偏向。但如果是一位趋势跟踪交易者，在整理行情产生的假信号中左右挨打了十次。那么第十一次信号的真假，真的与此前十次假信号没有任何相关性吗？

在我看来，对于寻找一套真正适合自己的交易系统而言，通过软件进行机械的统计，并且输出各种各样的统计数据的意义并不大。我比较倾向于将交易系统进行模拟交易，同时将战绩输出编制成一种系统指数。对于我来说，就编制出一种“一揽子期货品种趋势跟踪指数”。该指数就是假设在不使用资金杠杆的情况下，将资金平均分散在18个期货品种，并且用趋势跟踪系统进行模拟交易而生成的战绩输出。事实上，我就像研究上证指数一样，来研究我的那个“趋势跟踪指数”。渐渐地，我发现这种趋势跟踪指数的波动具有一些明显的规律。

我甚至希望通过对这些规律进行研究，从而对所谓的“趋势跟踪指数”进行交易。遗憾的是，用实战复制系统指数过于繁琐，并且很难在某一时点上对18个期货品种同时进行加减仓。另一方面，研究系统指数的减仓技术要比研究建仓技术困难得多，这或许是由系统的正期望收益所决定的。比如，2010年11月我的资金在12个交易日内出现了31％的资金回撤幅度，但这种资金大幅下跌只不过吞没了之前三周的涨幅罢了。也就是说，即使资金出现了大幅的回撤，但只有二十几个交易日的资金水平高于回撤后的水平！很显然，除非我能够精确地在这二十几个交易日内进行减仓，否则任何减仓策略都不会产生好的效果。而在此后8个月的时间内，我的资金都是在这12个交易日创下的高点与低点之间波动。大家是否能够体会到，试图逃避“交易的冬天”需要的技术水平苛刻到一种什么样的地步呢？与此对应的是，系统指数的走势图几乎从未跌破过60周（均）线，甚至历史上有很多次正好跌至60周线就开始大幅上涨。更有意思的是，历史上几乎没有一次顶部不是形成于从快速上涨到快速下跌的剧烈波动之中；同样几乎没有一次系统指数的牛市不是起步于下跌后形成的令人窒息的平静期。

虽然我无法直接交易系统指数，但我却可以将找到的某些规律运用于资金管理，甚至是调整在不同市场中的资金配置比例。此后我又编制了“工业品跟踪指数”“农产品跟踪指数”“股票指数跟踪指数”。但与此同时，我也发现了一个问题：资金大幅回撤后加仓是非常划算的，但此时也是资金最为匮乏的时候。从某种意义上来讲，就我的趋势跟踪系统而言，更高效的资金管理模式应该是“输冲赢缩”的。毕竟，快速亏损几乎完全来自于快速获利之后，而快速亏损的本身就是一种最好的风险释放过程。我突然发现，越容易发现的、越明显的规律，就越难以运用于实战获利。虽然说系统指数从未跌破过60周线，从而可以提示我在60周线处进行加仓。但历史上，有很多次“交易的冬天”在远高于60周线处结束。也就是说，如果我预留资金准备在60周线附近抄底的话，那么我很可能等到一次远高于现在点位的抄底机会。对于上涨而言，则更难以找到一种典型的上涨幅度，从而几乎让顶部减仓变得可望而不可及。

结论很有意思，如果考虑到系统的安全性，那么“赢冲输缩”就一定是对的；而如果考虑到系统的收益最大化，那么“输冲”就一定是对的。即使是在宏观上，系统也绝不可能保证一定会出现什么样的结果。历史上从未发生过的事情，并不代表未来不会发生。无论如何，对于一位稳健的交易者而言，都不应该不考虑破产风险而采用“输冲”的策略。最终我的资金管理被确定为“等额下注法”，也就是每次下注的手数都是固定的。当然，随着资金的增长，我总是要增加仓位。我更乐意于寻找正处于长期整理中的品种进行加仓。长期整理的品种必定会吸引大量的高抛低吸者参与，而他们买单之日就是我获利之时。从某种意义上来讲，在微观上我的确是采用“输冲”的策略，但在宏观上采用“赢冲”的策略。除非万不得已，否则我是不会进行可能丢掉位置的“缩仓”，当然这又让我不得不更加谨慎地应对加仓。

这或许提示大家，不要一味地人云亦云，而要更多地考虑自己的实际问题。“赢冲输缩”的价值在于最大程度地规避破产风险，同时又在最亏得起的时点上通过增加仓位来实现利润最大化。毫无疑问，“赢冲输缩”是有其巨大价值的。事实上，我的交易系统中，已经充分体现了“赢冲输缩”的思想。只不过我实现的方式并非是采用“资金投入量”，而是通过“资金利用效率”。

\subsection{分散的魅力}

{\kaishu
我相信，一套交易系统中，你决定的东西越少，市场决定的东西就越多。这样就能越接近市场的真相，并且产生更好的效果。

\hfill ——金融帝国
}

《通向金融王国的自由之路》中，将资金管理称为“可以使你在赢利的时候赚得更多，在亏损的时候输得更少”的技术。首先承认一点，我的数学水平不高，所以无法计算出具体的理论数值，那么我们就继续从概念的角度入手吧。毕竟精确的计算在混沌的市场面前，重要性将会大打折扣。我相信，沉重与压力同金钱无缘，那么今天的话题还是从游戏开始吧。让我们回到《通向金融王国的自由之路》中那个赌博游戏。如果你带着1000元参加这样一场赌博游戏，并且有100次的赌博机会，而且每次赌博的正确率为60％。你所能够左右的只是每次下注的多少。在你输的时候，会输掉你所下的赌注；而在你赢的时候，会赢得同你下注一样多的利润。

正如书中所说，输冲赢缩的下注策略，可以使赌客在拥有正期望收益的情况下最终亏损。那么什么是最佳的下注策略呢？按照书中的说法，每次用总资本的20％下注可以得到统计学上最优的结果。但作者并不推崇将这种下注策略运用到期货市场，毕竟它可能会产生超过80％的资金回撤幅度。这恐怕是所有期货交易者心理上所不能承受的。

回到我们熟悉的期货市场上来，问题可能会变得更复杂。毕竟我们下注所使用的保证金并不等同于我们可能会出现的亏损。同时，为了能够使交易系统按照固定规模运行下去，还需要留足结算准备金用于对潜在亏损的支付。好在我们无需过多地计算也能够得到一些常识性的结论：在具有正期望收益的期货交易系统当中，下注的多少可以与资金回撤幅度、收益率、破产风险呈现绝对的正相关性，并且这三者增加的幅度呈现递增的状态。看来多大的仓位才是最好的，还真的是一个无法回答的问题。

如果给上面的游戏，加上一个不能够使资金回撤超过30％的附加条件的话，那么可以通过计算得到每次下注金额的数值，我们先将这个数值称为X。这时我们将这个游戏稍微改变一下。其他的问题完全不变，我们把100次赌博变成100轮赌博。在其中每一轮的赌博中，我们都可以同时与10个不同的庄家进行赌博，并且概率和赔率完全不变。那么在保持相同的潜在资金回撤幅度的前提下，每轮下注金额的数值Y与此前的X会有什么样的不同呢？

即使我们没有任何数学知识，也不难得出Y＞X的结论。或者说，我们如果在每一轮下注中同样使用X的赌注，并且将赌注平均分散在这十个同时进行的赌博中，那么我们就可以得到更小的资金回撤幅度。如果我们要求的潜在资金回撤幅度不变，那么下更重的赌注就可以使得我们在100轮赌博后赚到更多的利润。当期望收益与交易频率都是确定值时，下注的多少就与获利的幅度呈现绝对的相关性。这就是分散投资的魅力！我们可以使资金波动得更加平稳，或者获得更多的利润。

这时出现了一个问题：上述游戏里每一轮的10次赌博中，只是期望收益是相同的，但赌博结果并不具有任何相关性。可是在复杂的金融市场上，我们很难说哪两个品种是完全不具有相关性的。当然，这种相关性即使在相同品种的不同月份上，也不可能是完全相关的。那么在真实的期货市场上，我们只能认为这种相关性介于0和1之间，而相关性的具体数值恐怕也是很模糊的。接下来，我们需要在考虑到不同品种相关性的情况下，来谈一谈关于分散投资的问题。

第一，指数具有策略统计的代表性吗？考虑到本轮大牛市我将更多的精力放在了期货市场，那么我们就用上一轮大牛市来举例吧。上一轮长期牛市如果用波浪理论来划分的话，那么牛市很可能开始于1998年8月的1043点，结束于2001年6月的2245点。我们假设有两位交易者，他们分别是A和B。其中，A对上涨指数ETF（当然这只是假设）进行跟踪；B将资金（按照权重）分散到几百支不同的个股上进行跟踪。并且他们采用完全相同的趋势跟踪系统，那么最终的战绩会有什么样的差异吗？

也许你并不相信，他们的差异是非常巨大的。B的战绩将明显好于A。很久以前我就发现了这种奇怪的现象：上证指数从1043点到2245点升幅为115％，而全部股票的平均升幅要大大高于这一数值。甚至好的股票可以上涨十几倍，而坏股票的升幅也很少出现低于115％的情况。这究竟是因为什么呢？也许你会对这个结论感到惊讶，那么就来翻阅一下股票的历史走势图吧。好在这是一个非常明显的现象，甚至不需要进行统计就能够被发现。

形成这种现象的原因，在于个股之间的相关性明显小于1。它们的最高点和最低点并非出现在相同的时点上，而这种差异还是非常巨大的。重组股的最低点在1997年底；低价大盘股的最低点在1998年8月；而老牌绩优股的最低点在1999年5月。从最高点来看：部分老牌绩优股的最高点在1999年7月；大部分重组股的最高点在2000年初。事实上，自2000年初亿安科技等科技股在疯狂中破灭以后，长达一年多的时间内，总会涌现出新的热点，以及最终见顶而开始下跌的股票。热点在科技股、低价股、钢铁水泥石化股、次新股之间转换不停。请注意，这种涨跌的差异性，在指数面前相互“中和”后彼此抵消了大部分的能量。

我在测试期货交易系统的时候，也在反复地研究期货指数与真实交易战绩之间的差异。期货指数同样存在“掩藏趋势”的特征。这是由于换月等原因出现的“误差”。根据近月的趋势感强于远月的特征，我们在不断换月的情况下，跟踪铜价长期牛市得到的实际利润要高于跟踪铜价指数的利润。

第二，消失了的风险。我们可以通过单独计算各种期货合约在趋势跟踪后所产生的利润，然后将其相加就是我们能够得到的总利润。也就是说，我们在走势一定的情况下，只要机械地按照系统操作就能够得到唯一确定的收益。当然，这是站在事后统计的基础上。如果我们回到过去的交易过程中，就会发现一种有趣的现象。

如果按照趋势跟踪系统不加思考地应对市场，那么我们就会经常出现“无效”持仓的情况。比如说，我现在持有玉米近期合约的空头，同时持有玉米远期合约的多头；持有小麦合约的空头，但却持有黄豆合约的多头；持有橡胶的多头，而又持有铜的空头。就目前的持仓状况来看，我也只不过是在工业品上偏空、农产品上偏多。而总体上我为空头头寸支付的保证金，几乎和多头头寸一样多。

很多人认为我是完美主义者。也许过去是，但现在真的不是了。以前我总是要么全对、要么全错。但现在我“多空对持”的做法，从根本上就注定了必然有一部分的持仓是“错”的。这种不可能“全对”的做法，恐怕很少有人愿意放弃主观判断来做到这一点。请注意，这种“多空对持”并非是有意进行套利。它既不是我的有意所为，也不是我有意避免的。我只不过在做市场告诉我的事情。

这种多空对持的无效持仓，根本不可能直接给我们带来利润，同时也自然不会直接带来风险。那么这种无效持仓究竟有没有意义呢？如果我们将这种“对持”拆成孤立的两个部分，那么他们应该是具有相同期望收益的不同交易。从宏观上来讲，我们并没有丧失任何统计中的利润，或者说整体利润就是两笔不同交易之和。前面已经说过，当期望收益与交易频率一定的时候，仓位增加一倍、利润就会对等地增加一倍。

可是在微观的时点上，这种能够带来长期宏观利润的风险消失了。我的意思是，这种与宏观利润对应的微观风险在分散投资形成的“多空对持”面前消失了，但这又并不妨碍我们得到宏观的利润。

当我们把两笔各自独立的交易所形成的资金波动合二为一时，显然会出现整体资金波动幅度的“抹平效应”。如果你的交易系统具有正的期望收益，那么宏观上的获利就存在必然性。而所谓的风险，只可能来自于微观上，这就是“抹平效应”的价值所在。毕竟，让我们所惧怕的破产风险，只可能来自于市场极端波动所形成的微观风险。请注意，如果某个交易系统具有宏观风险的话，那么就不是资金管理能够解决的问题了。

很多人觉得我的交易系统太过僵硬，如此僵硬的方法无法应对混沌的市场。事实上，真相绝对是截然相反的。我用僵硬不变的系统，采用僵硬不变的参数，跟踪僵硬不变的品种，甚至连下注的比例都是僵硬不变的。但请注意，我应对市场的风险水平却是在杂乱无章地、无迹可循地、大幅度地随着市场进行变动。我在任何一个时点上所承受的风险水平，都完全来自于市场的提示。我相信，一套交易系统中，你决定的东西越少，市场决定的东西就越多。这样就能越接近市场的真相，并且产生更好的效果。

想想看，我持有的众多品种在出现一致的持仓方向之前会先出现什么情况？那就是获利！多空略微的偏离，就能带来略微的利润。如果真的出现一致的持仓方向，那么我已经出现了不小的利润。请注意，潜在风险与潜在获利都在随着获利而增加，同时会随着亏损而减少。即使我的总体仓位是固定的，那么也可以通过资金利用率，让市场自己来调控“赢冲输缩”的程度。

第三，不安分的市场。运动是宇宙的本质，那么运动就是绝对的，而静止才是相对的。无论我们如何降低温度，都无法使分子真的安静下来。宏观上的平静与微观上的惊涛骇浪并不矛盾。当你看到一杯平静的水时，是否能够想到其中分子运动有多么激烈呢？

我们可以轻松地找到一个外表平静的市场，但我们却很难找到一个“从里到外”都平静的市场。虽然，我在期货市场上交易的时间还不长，但我也能感受到这种现象。我们无法否认，从2007年初到8月份的农产品市场处于一个明显的整理阶段。但在这种平稳的外表下，内在的部分又是如何呢？豆油大幅度地走高；玉米快速凶猛地暴跌；小麦5月合约不可思议地大幅下挫。当然，黄豆、豆粕、小麦远期合约的确进行窄幅整理。

对于股市来说，我有更长时间的观察经验。1998年前后长达两年的横向整理中，造就了垃圾股的疯狂，同时也使老牌绩优股大幅暴跌。这一点我的印象最为深刻，因为我当时持有的就是老牌绩优股。2002年至2004年的股市确实非常平静，但其中也上演了汽车股的疯狂，核心资产类股票的中期行情。当然，我持有的垃圾股照样大幅度下跌。

我们是否可以这样理解？市场整体的趋势，必然代表着更多个体的趋势；但市场整体的整理，并非代表全部个体的整理。即使是表现最为平淡的年份，无论是股票还是期货市场，都会涌现出黑马和黑熊。假设市场的整体走成静止的一条水平线，那么微观个体也是会存在毋庸置疑的差异性。很多人都在绞尽脑汁地解决选股问题，那么就说明个股之间的相关性的确不会太大。与此同时，个体之间的相关性越低，那么它们的波动差异性也就越大。我们再来“偷换”一下概念，将差异性视为趋势性。那么我们就可以得到一个有意思的结论：只要个体之间的差异性大于我们的趋势跟踪尺度，那么我们就可以在整体最为平静的时期获利！当然，品种间的相关性与趋势跟踪尺度是在不同角度下的产物，或许它们并不具备绝对的可比性。但我相信，它们之间或多或少会有一些关系。

回到我们的趋势跟踪系统上来，我们已经得到过这样的共识。趋势跟踪者不怕上涨，也不怕下跌，甚至连整理都不怕。小的整理能够被我们包容，大的整理反而能够让我们获利，我们怕的只不过是那种不大不小、恰到好处的整理。请注意，趋势跟踪系统能够呈现风险有限而利润无限的特征。也就是说，价格波动在一定区间之内能够亏损，而这个区间之外都能够获利。

均线形成的趋势跟踪系统无法清晰地量化，那么我们就“哈奇10％投资法则”来进行说明吧。这个法则很简单，就是从最高点下跌10％卖出，从最低点上涨10％买入，并且长期不间断地跟踪。假设价格在9元至10元的价格区间内不断地上下震荡，那么“哈奇10％投资法则”就会陷入一种泥潭。这就是那种“恰到好处”的整理，也是我们能够假设的最为糟糕的情况。如果价格波动幅度能够缩小一点，那么就能够被这种方法所包容。如果价格波动幅度能够扩大一点，那么就能够减少一点亏损。直到震荡区间扩大到20％，就会出现盈亏平衡点。很显然，这时的亏损区间为10％至20％的波动幅度。小于10％的波动幅度将会被系统所包容，而大于20％的波动幅度将会使我们实现利润。

再回到相关性的问题上来。假设某个品种的整理尺度正好同我们的跟踪尺度“恰到好处”，那么其他品种会出现什么样的情况呢？很明显，差异性就等于了偏离程度，并且使得其他品种的整理幅度不可能同样的“恰到好处”。这似乎就已经证明：当不同持仓品种之间的差异性能够大于趋势跟踪的亏损区间时，那么其他品种就不会出现亏损，甚至还会出现利润。

分散投资由于各品种间差异性的缘故，无疑将“哈奇10％投资法则”的理论最大亏损值减小了。也就是说，同时出现10％的亏损是不可能的！当然我也不得不承认，最大亏损值具体降低了多少，仍旧是我无法计算的。而我也只能从概念上认为：分散的品种越多，品种之间的相关性越小，那么最大亏损值就会降低得越明显。简单地说，当分散得越充分时，触及跟踪系统的死穴就会变得越困难。既然趋势跟踪系统只会去回答市场出现什么样的走势，才会出现最大的亏损。那么，对于最不利走势的要求就已经变得越来越苛刻了。没错，所谓恰到好处的定义，已经变得越来越复杂，甚至复杂到我都无法去计算。

这样一来，交易系统就能够在分散以后，出现集中操作不可能具备的稳定性。也许很多人立刻会反对，认为我的说法似乎太过于抽象。毕竟，假信号的数量并不会被这种分散所减少。没错，我绝对没有说分散投资能够消灭风险。我只是说，它能够在同一时点上把微观的风险降低。事实上，当你把资金分散得足够广泛时，就会在任何时点上都会遇到假信号的骚扰。没错，从来都不间断的骚扰！我也从来都没有说过，资金管理能够改变你的盈亏，这是由系统的期望收益来决定的。资金管理的魅力在于，它能够让你更平稳地赚钱或者亏损。至于为什么能够让你赚得更多，则是因为在足够平稳的情况下，你能够通过增加仓位来实现更大的利润。

如果你质疑趋势跟踪系统的价值，那么讨论这种资金管理就是没有意义的事情。这里我们假设，如果你能够永远用固定的资金来进行不间断趋势跟踪，那么足够长的时间后你会必然盈利。也许你已经明白了，资金管理的价值并非消除宏观上的风险，而只是抹平微观上的风险。事实上，分散的意义只不过是把你在不分散投资的情况下，短时间内面临的亏损，拉到更长的时间里来释放。再通俗一点讲，假设你某个月可能会亏损12万元，并且对交易系统带来严重的冲击，那么资金管理策略并不能消除这12万的亏损，但它却可以把亏损分散在一年的时间内消化。当然，这只是理论上的假设，而现实绝对不会如此完美。相信我，当你能够在期货市场上得到稳定性，那么你就可以变通地得到你想要的一切。阿基米德说过：给我一个支点我能撬起地球。而我觉得，给我足够多的非相关品种，那么我将能够把资金曲线拉成直线。当然把资金曲线拉成直线是不可能的事情，毕竟金融市场上严格意义的不相关品种是不存在的。甚至在极端行情面前，通常不相关的品种也会出现明显的联动，那么我们就绝对不能在确定仓位时忽略了谨慎性原则。

显而易见的是，太大的资金回撤幅度能够最大程度地阻碍合成复利的梦想。我并不想更多地改变自己使用的趋势跟踪系统，而希望更多地从资金管理入手。这样就可以尽量提高稳定性。其实，对稳定性的追求，无异于对宏观确定性的追求。我越来越相信“微宏互换”中所讲述的道理。当我放弃了微观上的确定性以后，却在如此狂热地追求宏观确定性——稳定性！

\subsection{对资金管理的思考}

{\kaishu
资金管理角度上的策略，只可能对缺陷进行转化，而不可能直接将缺陷消灭。

\hfill ——金融帝国
}

我对于趋势跟踪的了解可以追溯到十多年前，而对于资金管理的了解则是最近几年进入期货市场之后的事情。从某种意义上来讲，资金管理的意义就在于规避破产风险，而股票市场几乎并不存在破产风险。事实上，更多股民对于资金管理的“研究”只不过是为了取悦自己的心灵罢了。而我则曾经戏称：“牛市继续的话，我留钱有何用？牛市完结的话，我留钱又有何用？”我的意思是，股价下跌以后我也不会抄底，相反会卖出更多的股票。这时我不但不需要资金，反而会使资金更加富余。即使牛市完结了，我也不过是付出一次止损成本罢了。

我不知道大家能否接受，我这种对资金管理“东拉西扯”式的讨论方式？无论大家是否接受，或许我也只能够用这种方式来与大家讨论了。记得上中学的时候，我的物理成绩非常好，甚至物理竞赛能够拿到全年级第一名。有意思的是，我一看物理公式就脑袋发蒙，但只要看一遍书就能理解其中的原理，从而能举一反三解答难题。或许，我从来就不喜欢用公式般的逻辑来思考问题，而是通过理解概念的方式来看待问题。很显然，我无法细腻地回答出任何关于资金管理的问题，更不可能写出复杂的公式来应对这些问题。当然这又未必不是这本书的一个优势，毕竟形象的东西更容易理解。既然本章节的题目叫做“对资金管理的思考”，那么我就不追求得出什么严谨的结论，而是采用论坛式的讨论风格。

前面已经说过，资金管理策略中并非存在某种最佳策略，而是一种非常个性化的技术。在对交易系统的研究过程中，我发现了大量的“互损原则”。我们在某个环节得到一些什么的时候，就注定要在另外的环节失去一些什么。也许真的像《孙子兵法》说的那样：故备前则后寡，备后则前寡；备左则右寡，备右则左寡；无所不备，则无所不寡。其实，资金管理的意义在于将利弊进行能够让交易者接受的方式转化。从某种意义上来讲，脱离了行情系统来探讨资金管理策略是没有任何意义的。相对于学习而言，我更乐意于思考。这恐怕是这本书最为明显的一种风格：我会站在自己的角度来对正统的权威观点进行思考，而思考的本身就是一种质疑的过程。何况，对于更加个性化的资金管理策略，似乎我们就更不应该无条件地接受正统的观点。

比如，我是通过“占用保证金的数量”来确定单品种的仓位，而不是通过“真实波动幅度”（N值）来确定单品种的仓位。当我身边的朋友想学习炒股时，我推荐他看一看《海龟交易法则》这本书。而看完此书后，他认为正确计算N值才是交易获利的关键。要知道这是一位从来没有接触过股票的初学者，而他感兴趣的居然是本书中最为“复杂”的部分，毕竟这是书中唯一需要计算的内容。当然，这或许与他是软件工程师有关，但这种倾向似乎也在说明交易者更希望通过类似于科学的方式来应对市场。事实上，期货合约的保证金率，本来就是根据过去的波动幅度，而对潜在风险的一种“预测”。我就是这么懒惰，有捷径就不想再去费事，而不是用最精确的方法来应对市场。虽然保证金率与N值都来自于过去价格的波动幅度，从而无法准确反映出未来可能突然增加的波动性。也就是说，他们同样对于风险的表述具有明显的滞后性。但保证金率似乎考虑了更多的问题，从而使它远没有N值那样敏感。比如说，股指期货的保证金率大大高于商品期货，这并非完全因为波动幅度的问题，而是考虑到涨跌停板幅度等交易规则。事实上，随着期货品种从平静走向热络，N值有可能增长数倍，但保证金率的提高往往很难超过一倍。

如果按照“真实波动幅度”来确定仓位的话，那么我就会在2008年3月的棉花和铝合约上建立一个令人惊讶的大仓位！毕竟，棉花在“多头陷阱”前的波动幅度是非常小的。但在“多头陷阱”中，棉花的波动幅度丝毫不逊色于波动幅度最大的品种。相反，它在高位停留了更短的时间，并且成为我亏损幅度最大的一个品种！要知道，无论你是用均线死叉，还是用跌破n日内的最低点作为离市信号，在顶部停留时间短的品种都会给你带来更大的亏损。风险来自于微观上的极端波动，而我更情愿相信所有品种最不利的极端波动都是一样的！对于2008年10月的全球金融危机来说，无论你持有什么品种的多头仓位，恐怕都免不了品尝跌停板的滋味。

当然，我也并非是在否认“真实波动幅度”的价值。我猜测海龟交易法则的本意在于：历史上很多壮观的大趋势都起源于极度平淡的市场阶段，从而使得交易者因极低的N值而建立一个巨大的头寸。很显然，趋势特征越明显，N值就会越大；那么趋势特征越明显，对于建仓而言就应该越谨慎。前面我已经说过，市场的趋势特征越明显，我们就越难以控制风险，毕竟“证实”趋势的完结需要更大的价格折返幅度。事实上，我的交易系统相对于海龟系统而言更加苛刻。甚至不会在趋势明朗之时建立仓位，而是通过在整理行情中的“试错”来捕捉未来潜在的趋势。通俗地说，本人进行连续不间断跟踪的开端，一定是N值很小的整理阶段。这里我想说的是，对于实现某种交易思想而言，完全可以通过截然不同的方式来实现。当交易者借鉴交易大师所提倡的技术时，如果并不能深刻理解这种技术的本意，那么生搬硬套的借鉴也是没有任何意义的。我们应该更多地考虑技术之间的兼容性，而不是绝对忠于技术本身。当我完全避免介入微观上N值巨大的市场阶段时，那么宏观上波动性巨大的品种反而具有更大的优势。当我假设“任何品种最不利的极端波动都是一样的”，那么波动性大的品种，潜在利润与潜在风险的比值就会越大。

这里我来详细地解释一下，我是如何规避趋势特征明显时的建仓行为。表面上来看，趋势跟踪系统是一个“求势而舍价”的思维方式，但实际上并非如此。其实捕获大趋势是非常不容易的一件事情，我们必须具备趋势位置的优势，同样也必须具备价格位置的优势。如果你没有一个良好的入场价位，那么你会更容易丢掉趋势中的位置。只有一个良好的入场价位，你才能够心平气和地忍受价格的正常回撤。

即使一只股票从10元上涨到100元，如果我从11元买入而你在20元买入的话，那么最终的结果很可能并非你只少赚了9元。20元买入的你，别说遇到轻微的回调就会仓惶而逃，甚至继续上涨都能给你带来一种恐惧。趋势感越强的行情，往往证明趋势完结就需要越大的折返幅度。也许该股从20元跌到15元才能“证明”趋势的完结，而你恐怕无法承担这5元的价格折返，但11元买入的我却可以欣然接受。

这就是我坚信的一个原则：永不非信号开仓！如果我的每次开仓都对应一个信号的话，那么我的潜在亏损也必然符合跟踪系统的亏损分布。但如果我开仓点位没有对应一个信号，那么我很可能在承受巨大的资金回撤与丢掉位置之间左右为难。从某种意义上来讲，我们的入场信号就是在趋势与价位之间找到的一个平衡点。趋势与价位是互损的，同时并非趋势感越强越好！这就是盈利加仓的一个障碍。我无法想象，在2008年3月的资金高点时，我进行加仓会有什么样的后果。也许，我这样做反而会在那次“多头陷阱”中亏损更少的钱，毕竟不好的价格位置更容易让我成为惊弓之鸟。但我无疑承担了巨大的丢掉位置的风险。

连续不间断趋势跟踪还有一个非常重要的要求，那就是必须产生“绝对信号”而不会让我们错过任何趋势行情。即使你的开仓都会对应一个入市信号，但如果你不连续跟踪也可能产生不良的结果。简单地说，假设我的趋势跟踪系统具有三分之一的正确率，那么有两种不同的应对方式。我在一个有2个绿色球和1个红色球的袋子中，连续三次拿出一个球，那么两个绿球让我亏损两次，而一个红球让我获利一次。但如果你的袋子里有2000个绿球和1000个红球，你连续取出一些球并承担盈亏的话，即使我们的概率都是相同的，但你无疑面对比我更大的破产风险。因为我有固定的利润对冲周期，而你却并没有！

连续不间断跟踪其实还能够通过机械的方式，有效地控制好交易节奏。我们假设有十二个袋子，每个月只能在一个袋子中抽取球来决定盈亏。能够让你获利的红球整体上仍旧只有三分之一，但其中有4个袋子中有大量的红球，而其中8个袋子中有大量的绿球。虽然抽取红球所对应的利润要大大高于绿球所对应的亏损，但我们却只有在抽取到绿球时才有继续抽取的权利。当我们抽取到红球时，就只能等到下个月再从其他的袋子中继续这个抽球游戏，毕竟红球的利润来自于长时间的持仓。如果我们并不知道每一个袋子里球的颜色分布，那么每个袋子只取1个球恐怕就是最为明智的选择。但对于不连续跟踪的交易者来说，他们很可能在几乎全是绿球的袋子中反复不断取出球来。遗憾的是，我见到太多不懂得休息的股民，在2008年的大熊市中不断寻找买入信号建仓，而结果只是一次又一次止损离场。学会炒股获利并不困难，只不过需要学会三件事情：敢赢（敢于长时间持有获利筹码）、敢输（敢于果断的进行止损）、敢等（在没有获利机会的市场阶段什么都不做）。要知道，学习这三件事情的难度是逐级递增的，而学会等待恐怕是最为困难的事情。对交易没有热情的人，往往难以付出足够的精力来学习交易；而对交易有足够热情的人，往往更加难以学会等待。

能够影响到我们账户存亡的风险只有两种：第一，连续不断出现的小亏损；第二，价格突然大幅折返的大亏损。在我看来，后者通常是毁灭账户最大的潜在风险。我相信，正统的资金管理法则是可以保护投资者的安全的。比如，《走进我的交易室》介绍的“2％、6％”法则就是一个不错的主意。我们在任何一个仓位上都不能承担超过2％的初始风险；而任何时候我们都不能承担超过6％的总体初始风险。但我觉得，这其中也有一些明显的弊端。

我们假设期货的保证金率为10％，考虑到隔夜跳空、交易成本、滑移差价等现实因素后，过于紧密的止损点就是毫无意义的。如果我们按照4％的止损幅度来计算2％的暴露风险所支持的仓位，那么我们只能够建立一个只有5％的仓位。也就是说，我们在6％的总体暴露风险的基础上只能够建立一个15％的仓位！很少有人甘心接受如此轻的仓位，那么就有可能出现降低止损幅度的误区。既然规则是需要绝对遵守的，那么我们就只能变通一下，主观地追求更为紧密的止损点。我们可以把止损幅度降低为2％，这时我们就可以建立一个30％的仓位水平。我觉得30％的仓位应该是比较明智的，也是风险和利润的一个比较理性的平衡点。但问题是，2％的止损幅度很可能已经低于中期趋势跟踪敏感度的下限，单单是隔夜跳空就有可能高于止损幅度。这样一来，系统的稳定性以及可操作性就会大大降低。

其实，这与我当年研究的一种资金管理策略非常类似。我当时希望在每一次战略性建仓上承担10％的风险，而只有当这笔交易没有暴露风险后，才有权利进行下一步的建仓行为。当时我的止损幅度通常在4％左右，也就是可以建立一个25％的仓位。当我连续做对四次，那么我就能够将初始仓位打满。我觉得这种策略是可行的，但仍旧有一个明显的弊端。

事实上，这种风格体现了“赢冲输缩”的理念，但是与此同时却丧失了“择强汰弱”的优势。请注意，我们进一步建仓必须建立在我们此前交易获利的基础上。对于一个能够让我们获取暴利、具有足够市场广度的趋势行情来说，这种策略无疑让我们反应得太慢了。如果说我第一笔建仓已经没有任何暴露风险的话，那么我继续建仓的权利只有两种可能性：第一，找更为疲弱品种的买入信号来建仓；第二，在强势品种上非信号开仓。

如果一个品种发出一个明显的开仓信号，而我又没有开仓权力时，那么我只有放弃这个信号，从而丧失了连续不间断跟踪的可能性。也就是说，这种资金管理策略过多地考虑了交易者自身的问题，而同时过多地忽略了市场的问题！这时行情系统与资金管理策略是两个分别独立的系统。往往对于资金管理最适合加仓的时刻，恰恰是对于行情系统最不适合加仓的时刻。毕竟，你获利越多的时候就越有资格承担更大的风险，但这时的市场往往处于明显的趋势当中，从而使我们更加难以实现低风险的入市。按照《海龟交易法则》中讲述的一个故事，某人大言不惭地声称海龟的系统并不能获利，当然他并没有考虑到“择强汰弱”的法则。而海龟只会在一个板块中，选择最强的一两个品种建立多头仓位。似乎这种资金管理策略就注定了：你要么在最强的品种上获利后，在最弱的品种上建立多头仓位；你要么就在最强的品种上获利后，在同样强劲的品种上建立一个无法合理设置低风险止损点的多头仓位。在我看来，应该是先有交易系统，再有资金管理的问题。那么就应该永远让资金管理迁就交易系统，而不是相反。事实上，我更乐意通过交易系统的本身来实现资金管理的目的。与此同时，越简单的资金管理策略，就越容易产生好的效果。正如前面所说，我们应该让市场回答更多的问题，而让系统参数回答更少的问题。

其实，这种资金管理策略只能够让我们大体上保住了初始资本的安全，而同样无法避免大幅度的资金回撤。对于初始资本非常安全的保护，无疑大大丧失了获利的可能性！获利加仓的策略，使我们的盈亏平衡点大大提高了。也许只要行情出现一定程度的回撤，我们相对于初始仓位的“巨大”仓位就会轻松触及盈亏平衡点！当然，可取之处仍旧有两点：第一，你可以保护好你的初始资本；第二，你可以在疯狂的大牛市中建立一个超过初始资金100％的庞大仓位。这样不但可以获得暴利，而且你的初始风险几乎可以低到忽略不计。这两个优势非常诱人，但这是以缩减交易盈利区间作为代价的。绝对不可能亏损稍大一点资金，加上有获得暴利的可能，无疑会使你更加难以经常实现获利。更多的情况是，在一次次我可以获利的行情中，而你却一次次的“白玩”。当然，也许特定的行情，你可以赚取极其庞大的利润。这也许就是为什么交易大师说，市场中只有15％的时候是有趋势的。也许这种策略只有在15％的时候能够获取暴利，而其他时期都无法获利。

请注意，你总会面对一次最大的亏损！对于2008年3月“多头陷阱”的资金大幅回撤来说，我最终亏损离市的只有糖、棉花、豆粕、小麦这四个品种；而铝、铜、豆、橡胶都是获利离市的。要知道我在黄豆上的多头仓位是在2007年8月份建立的。这说明如果不考虑浮动盈利而只考虑初始风险的话，那么即使我的资金回撤40％多也并没有出现亏损。我始终不晓得为什么这么多人如此忌讳初始资金的亏损呢？这是否与大多数交易者忌讳亏损的原因别无二致呢？事实上，我早已将比十多年前投入市场的初始资金更多的现金拿到市场之外消费掉了。这样一来，我又该如何定义初始亏损呢？对于我来说，我只会关心一个问题，那就是资金回撤幅度。毕竟，浮动盈利也是真实存在的钞票，不是吗？

其实，我并非是在批判交易书籍上的资金管理策略，只是觉得孤立地讨论资金管理的问题没有任何意义。很显然，从对所有行情系统的适应性的角度来说，“2％、6％”的策略无疑是最好的。但对于交易者来说，更加需要的是具有针对性的良方，而不是“什么都治”的营养品。但反过来说，只要我们的行情系统之间存在一定的差异，那么我的美酒就很可能成为你的毒药。事实上，我的初始仓位在40％到50％之间，甚至在仓位达到75％以前都会坚持连续不间断跟踪而不缩减仓位。问题是，我能够使用如此“激进”的资金管理策略是建立在大面积分散投资、并且具有盈亏高度相关性的对冲机制的基础上。很显然，如果你与我的行情系统有一定程度的出入，那么如此重的仓位简直是一种疯狂。我们在这里讨论资金管理的意义也并非是告诉你一种可行的方法，而是给你一种分析问题的思路。

我在《分散的魅力》中讲述了一种市场自调节的风险控制系统。这或许给了我们一个有效的思路：从整体的角度，通过资金利用效率来实现资金管理的目的，而不是从个体的角度，僵硬地通过仓位控制来实现资金管理的目的。很显然，这是一个在整理市道中非常出色的风险管理策略，从而避免了“整理泥潭”中连续不断出现的小亏损所带来的破产风险。这种风险自调节系统在没有参数的情况下，靠一部分无效的不完美，来必然捕捉到另一部分有效的完美。并且充分体现了三种可贵的交易思想：顺势而为、择强太弱、赢冲输缩。

但这时又出现了一个新问题：虽然这种方法可以在并不具备明显“市场广度”的趋势阶段中捕捉到一部分品种的多头利润，或者另一部分品种的空头利润，但在最为强烈的趋势阶段中，不但会出现巨大的价格波动，甚至还会出现所有品种同方向运行的市场广度。随着利润极度膨胀，几乎全部品种不但具有持仓方向上的一致性，同时具有巨大的跟踪乖离率。很显然，我们将会面对巨大的潜在风险。也就是说，我们依然无法解决《走进我的交易室》中所说的那种价格突然大幅折回的大亏损。

当然，“2％、6％”法则同样无法解决这个问题。或许我们应该欣然接受这种大幅获利后出现快速亏损的缺陷，就是问题的解决之道。要知道，站在资金管理角度上的策略，只可能对缺陷进行转化，而不可能直接将缺陷消灭。难道我们真的应该接受这种缺陷吗？或者说，靠丧失1元的利润，来规避1元的亏损真的没有意义吗？这里需要解释一下，后面探讨的方法我并没有运用于实战，而只是一种理论上的设想。这或许是因为我受不了过于繁琐的测试，而在实战上的试验又需要漫长的时间才能够得出结论。

作为日线级别交易者的我来说，是不会在盘中进行任何决策与交易的，而是将每日收盘价作为最小的交易时点。即使是这样，对于不能依赖软件进行测试的我来说，通过手工来测试各个品种的模拟战绩曲线都将是一件非常困难的事情。为了简化测试流程，我发现可以将跟踪尺度的“敏感当量”用周线来表示。当我用5周与10周均线交叉作为信号来进行测试时，所生成的模拟资金走势图与我的真实资金走势图几乎是完全一样的。考虑到“本体相似性”的问题，我猜测小时图、日线图、周线图之间的波动规律应该不会存在任何本质性的区别。那么将我的交易系统应用于小时图，就理应同样能够产生获利的效果。而这种获利的效果，应该同用5日与10日均线交叉作为信号进行日线级交易别无二致。

这时就出现了一种“敏感系统”，而理论上它们应该与现有的“迟钝系统”具有相似的期望收益。只不过过于敏感的系统无法有效地包容隔夜跳空风险，甚至更快的交易频率还会使佣金与滑移差价快速累积。当然，更快的交易频率或许可以在相同期望收益下实现更大的收益率，但这并非是我关心的关键问题。我相信敏感系统并不具备相对于迟钝系统的交易价值，但却应该具备毋庸置疑的分散价值。既然交易品种的分散可以解决“连续出现小亏损”的问题，那么系统的分散是否可以解决“突然出现大亏损”的问题呢？

前面说过，对于趋势跟踪者而言，大的整理能够让我们获利，小的整理能够被我们包容，我们怕的是那种不大不小、恰到好处的整理。我发现一个有意思的结论：能够让敏感系统左右挨打的行情，的确可以被迟钝系统所包容；而能够让迟钝系统左右挨打的行情，的确可以使敏感系统获利。很显然，敏感系统的确具有分散价值。当然在整体市场处于整理阶段的时候，多空对持就能够使资金权益得到稳定性，而根本没有必要采用敏感系统。另外我们也不难发现：无论是敏感系统还是迟钝系统都具有相同的利润回收期。当市场出现明显的趋势行情时，无论你使用什么样的参数来跟踪趋势都可以实现利润。考虑到敏感系统的缺陷，我们完全没有必要在整理阶段采用系统间的分散。事实上，我对现有系统在整理阶段的稳定性非常满意。既然没有问题需要解决，那就完全没有修理系统的必要。而我们需要解决的问题，只是来自于具有明显市场广度的趋势阶段。在一致性的整体趋势行情中，各个品种的差异性突然消失了。它们更容易齐涨共跌，从而使我们根本无法得到分散投资的保护。

其实，我们需要研究的只是敏感系统在趋势行情中的表现。用越敏感的系统跟踪趋势行情，就会在距离最高点越近的价位实现逃顶。但与此同时，趋势中出现的无效假信号也就会越多，甚至多到我们根本无法获利的地步。我不会愚蠢到只因为敏感系统能够卖在更高的点位就使用敏感系统，而是会从整体来考虑此前众多假信号的成本。

我们不妨设想这样一套系统：当市场出现一致性的上涨时，只要资金曲线的背离率大到一定程度，就将所有持仓品种切换到敏感系统。也就是说，行情的下跌触发敏感系统的信号，我就反手做空。与此同时，只要有三分之一的品种已经持有空头头寸，那么剩下三分之二的品种就继续切换回迟钝系统。如果未来行情最终证实下跌只不过是调整，那么翻空的品种就继续通过敏感系统实现翻多。随着时间的推移，资金曲线的背离率开始缩小，当其小于切换敏感系统的数值后，那么全部品种就继续使用迟钝系统进行跟踪。乖离率的缩小等于潜在风险的降低，同时这可以避免敏感系统在趋势行情中陷入多次小幅上下整理的泥潭。

我们不难想象，只要趋势行情还在继续，那么使用敏感系统就是一种增加微观亏损的行为。好在我们只不过是在最弱的三分之一品种上做空，那么结果只不过是降低了获利的速度，而绝对不可能出现宏观上的亏损。我们通过降低获利的速度，规避了“承担大幅资金回撤与放走行情”之间的矛盾。

当市场最终形成顶部时，敏感系统可以对资金快速回撤起到明显的减震作用。当然我必须承认，即使在顶部形成之时，也是可以出现引入敏感系统无法产生正面效果的可能性。毕竟顶部在形态上也可能是多种多样的，而不一定是我们所惧怕的那种尖尖的、快速下跌的顶部。即使这样我们仍旧不能认为敏感系统是毫无作用的。我们一直在讨论系统的“命门”，也就是市场如何运行才会使我们的系统出现最不利的情况。事实上，越复杂的系统分散就越难以出现，市场真实走势恰巧与“系统最不利的情况有”完全一致的可能性。至少，这种最不利情况出现时的可怕程度将会大大的降低。我们很难想象，当各个品种在出现巨大乖离率的情况下快速下跌，并且在深幅下跌使得持仓大面积翻空后，又出现大幅反弹的情况下，会给资金权益带来多么可怕的打击。除非我们完全没有利用浮盈加仓，否则后果真的不堪设想！当然，这种可怕场景真实出现的概率在理论上是非常小的。

这里我们没有必要说得太过复杂，一个简单的比喻就能够说明问题。我们只不过是在获利的黄金阶段拿出一小部分利润购买保险，然后用做资金大幅回撤时期的减震器。很显然，在绝对不能丢掉趋势行情中位置的前提下，规避资金大幅回撤的方法只可能是在此之前降低获利速度。那么这种靠放弃更多利润来规避亏损的转化真的有意义吗？

从某种意义上来讲，系统的分散与品种之间的分散别无二致。多空对持的确可以更大程度地体现出概率优势的魅力，那么刻意地追求多空对持就是一种折衷的选择。与此同时，两种不同形式的多空对持，同样可以体现择强汰弱与赢冲输缩的优势。事实上，资金管理全部的意义只不过是：通过各种各样的方法将盈利阶段的利润用做抵消亏损阶段的损失。很显然，敏感系统并不能增加期望收益，相反有可能降低收益水平。否则话，我就不会只是将敏感系统作为应急系统了。或许这种资金管理策略有可能是为了规避1元的损失，而放弃了1.2元的利润。我们可以换一个角度来更加抽象地分析一下这种转换的价值。所谓期望收益就是我们承担1元的潜在风险，能够获得x元的潜在收益。假设，我们的期望收益为2元。这就是说明，我们面对了1元的潜在风险，以及2元的潜在利润。微妙之处在于，如果我们能够把未来的潜在利润“借来”1元，来对冲掉1元的潜在风险，那么这种策略没有价值吗？事实上，“借来”的1元潜在利润，恐怕只能对冲掉0.8元的潜在风险。这时的期望收益成为了0.2元的潜在风险，对应1元的潜在收益。也就是说，期望收益已经变为了5，比此前的2大大提高了。当然，增加利润的方式来自于仓位提高了1.5倍，而我们微观上的风险水平并没有提高。

在微观上“借来”并不属于你的东西，就可以在不违反“能量守恒定律”的基础上，使你在宏观上得到不可思议的效果！量子理论中的“隧道效应”指出，电子可以“借来”并不属于它的能量，而使它穿越壁垒后再偿还这部分能量。如果说，电子可以在不违反能量守恒定律的情况下穿越壁垒，那么，我们也可以在行情系统期望收益不变的情况下获得更多的利润。我想告诉大家的是：对于任何不同的交易技术而言，即使不可能出现绝对的“利弊守恒”，也会表现出明显的“利弊相关”。对于我的思维方式来说，从来都不会奢望能够彻底将弊端消灭。遗憾的是，更多的交易者则热衷于将弊端隐藏起来，却自认为已经将弊端彻底消灭了。他们会使用牛市做多、熊市做空；整理阶段使用摆动类指标、趋势阶段使用趋向类指标；复杂顶部使用迟钝系统、简单顶部使用敏感系统等，以此来解决现实中的实际问题。

\clearpage
